\exercisesection{}

\begin{enumerate}
\def\labelenumi{\arabic{enumi})}
\tightlist
\item
  设 X 为局部紧空间，\(\mathfrak{F}\) （或 \(\mathfrak{F}(X)\) ）为 X 的闭子集的集合。对 X 的每个紧子集 K、K 的每个紧邻域 L 以及 L 的唯一一致结构的每个邻偶 V，令 \(Q(K,L,V)\) 为 \(\mathfrak{F}\) 中满足下列两个条件的元素对 (A,B) 的集合
\end{enumerate}

\[
\mathrm{A} \cap \mathrm{K} \subset \mathrm{V} (\mathrm{B} \cap \mathrm{L}) \quad \text { and } \quad \mathrm{B} \cap \mathrm{K} \subset \mathrm{V} (\mathrm{A} \cap \mathrm{L}).
\]

\begin{enumerate}
\def\labelenumi{\alph{enumi})}
\item
  证明诸集合 \(\mathbf{Q}(\mathbf{K},\mathbf{L},\mathbf{V})\) 构成 \(\mathfrak{F}\) 上一个分离一致结构的邻偶基本系。
\item
  设 \(\mathcal{U}\) 是一个与 X 的拓扑相容的一致结构。对 X 的每个紧子集 K 与 \(\mathcal{U}\) 的每个邻偶 W，令 P(K, W) 为 \(\mathfrak{F}\) 中满足下列两个条件的元素对 (A, B) 的集合
\end{enumerate}

\[
\mathrm{A} \cap \mathrm{K} \subset \mathrm{W} (\mathrm{B}) \quad \text { and } \quad \mathrm{B} \cap \mathrm{K} \subset \mathrm{W} (\mathrm{A}).
\]

证明诸集合 \(\mathbf{P}(\mathbf{K},\mathbf{W})\) 构成 \(a\) ) 中定义的一致结构的邻偶基本系。

\begin{enumerate}
\def\labelenumi{\alph{enumi})}
\setcounter{enumi}{2}
\tightlist
\item
  设 \(X'\) 是 X 的 Alexandroff 紧化，\(\omega\) 是 \(X'\) 的无穷远点；对 X 的每个闭子集 A，令 \(A' = A \cup \{\omega\}\) 。证明映射 \(A \mapsto A'\) 是 a) 中定义的一致空间 \(\mathfrak{F}(X)\) 到 \(\mathfrak{F}(X')\) 的由含 \(\omega\) 的闭子集构成的子空间上的一个同构（利用 b)）。由此推出 \(\mathfrak{F}(X)\) 是紧的（cf.~GT, II, §4, Exer. 15）。
\end{enumerate}

\begin{enumerate}
\def\labelenumi{\arabic{enumi})}
\setcounter{enumi}{1}
\item
  设 \(G\) 为局部紧群，\(\mathfrak{F}(G)\) 为 \(G\) 的闭子集的一致空间（在 No.~6 或 Exer. 1 中定义）。直接证明闭子群的集合 \(\Sigma\) （\(G\) 的）在 \(\mathfrak{F}(G)\) 中是闭的（不利用 No.~1 的命题 3）。
\item
  设 \(G\) 为局部紧群，\(\chi\) 是 \(G\) 在 \(\mathbf{R}_+^*\) 中的一个连续表示。把命题 4、5 与 6 推广到集合 \(\Gamma_{\chi}\) （诸测度 \(\alpha \in \Gamma\) 所成的），其中 \(\chi\) 在 \(H_{\alpha}\) 上的限制是 \(\alpha\) 的模。（把测度 \(\mu\) 换成 \(\chi \cdot \mu\) ，并考虑商测度 \((\chi \cdot \mu) / \overset{\vee}{\alpha}\) 。）
\item
  设 G 为一个不能由 G 的任何紧子集生成的局部紧群，并设 H 是 G 的一个使 G/H 为紧的离散子群。证明 H 不能由有限多个元素生成，但 H 的规范化 Haar 测度 \(\alpha_{0}\) 属于 \(\Gamma^{0}\) 中诸 \(\alpha\) （满足 \(\|\mu_{\alpha}\| = +\infty\) 者）所成集合的闭包（考虑 H 的由有限多个元素生成的诸子群的规范化 Haar 测度）。
\item
  设 G 是一个作为二元素子群的无穷序列 \((\mathrm{G}_{n})\) 的直和的离散群，并设 \(p_{n}\) 是 G 到 \(G_{n}\) 上的投影。令 \(\mathrm{H}_{n} = \overline{p_{n}}(e)\) ，并设 \(\alpha_{n}\) 是 \(H_{n}\) 的规范化 Haar 测度。若 \(\alpha\) 是 G 的规范化 Haar 测度，证明在 \(\Gamma_{c}^{0}\) 中序列 \((\alpha_{n})\) 收敛于 \(\alpha\) ，但 \(\|\mu_{\alpha_{n}}\| = 2\) 且 \(\|\mu_{\alpha}\| = 1\) 。
\item
  设 G 为不满足 No.~4 的条件 (L) 的局部紧群。
\end{enumerate}

\begin{enumerate}
\def\labelenumi{\alph{enumi})}
\item
  证明（用 No.~4 的记号）：N 到 D 上的典范双射不是同胚。（对 G 中 e 的每个邻域 W，令 A(W) 为含于 W 中的有限子群 \(\neq\{e\}\) 所成的集合；证明诸 A(W) 构成一个在 D 中收敛于子群 \(\{e\}\) 的滤子的基，但若 \(\alpha_{H}\) 记 H 的规范化 Haar 测度，则诸测度 \(\alpha_{H}\) 不收敛于 \(\varepsilon_{e}\) 。）
\item
  此外若 G 是可度量的，证明 G 的离散子群空间 D 中存在一个紧子集 B，它不含于 No.~3 的 Th. 1 中定义的任何集合 \(D_{U}\) 之中。
\end{enumerate}

